\scanpage{317}

\mspage{56}
% archive sheet number: 316

d'une suite de $n$ générateurs parmi les $f_{\varepsilon,\ell}$. L'application de dégénérescence $\operatorname{Nerf}_1\unsure{\text{ --- }|\longrightarrow}\operatorname{Nerf}_2$ est \unsure{ainsi|ici} donnée sur les générateurs par
\[
f_{\varepsilon,\ell}\longmapsto f_{\varepsilon,1,\ell}.
\]

Heureusement, on n'a plus à prouver l'associativité de la composition.

\uline{NB} La simplicité de ces descriptions \deleted{\unsure{me}} renforce mon sentiment que la catégorie $\underline{\mathrm{Ty}}$ définie à présent, avec les \uline{morphismes d'intervalles orientés croissants} $I_\tau\text{ --- }I_{\tau'}$ comme flèches de $\underline{\mathrm{Ty}}$, est bien ``la bonne''. On va déjà voir si elle donne une description sympa des catégories de fractions $\unsure{X.\Sigma^{-1}|X\Sigma^{-1}|Y.\Sigma^{-1}|Y\Sigma^{-1}}$, sans tortillements comme précédemment!

Mais on a promis de déterminer $\operatorname{Nerf}_r(\underline{\mathrm{Ty}})$, pour \unsure{tt|$\forall$} $r\in\mathbb N$. \deleted{L'élément \unsure{gé}} Le générateur général de $\operatorname{Nerf}_r$ est de la forme (\inserted{pour} $\varepsilon\in\{\pm1\}$, fixé)
\diagram{\[
\begin{tikzcd}[ampersand replacement=\&,column sep=small]
\underset{\substack{\Vert\\\Delta_\varepsilon^{\ell_0}\ \text{avec }\ell_0=1}}{\Delta_\varepsilon}\arrow[r,"f_1"] \&
\Delta_\varepsilon^{\ell_1}\arrow[r,"f_2"] \& \Delta_\varepsilon^{\ell_2}\arrow[r] \&
\cdots\arrow[r,dash,"f_r"] \& \Delta_\varepsilon^{\ell_r}
\end{tikzcd}
\tag*{\unsure{($\alpha$)|(3)|ds}}
\]}
où les $f_i$\inserted{$(i\in[1,r]):\Delta_\varepsilon^{\ell_{i-1}}\longrightarrow\Delta_\varepsilon^{\unsure{\ell_i|\ell_{i-1}}}$} sont donnés par \deleted{\unsure{des}} \deleted{\unsure{p}} applications
\[
\pi_{i-1}:[1,\ell_{i-1}]\to\mathbb N\qquad
\text{telle que}\quad
\underset{\substack{\Vert\,\mathrm{déf}\\
\sum_{j\in[1,\ell_i}\pi_i(j)}}{\operatorname{long}(\pi_{i-1})}
=\ell_{\unsure{i|i-1}}
\]
(N.B. $\pi_0:[1,1]=\{1\}\to\mathbb N$ est \unsure{identifié|identifiée} à $\operatorname{long}(\pi_0)=\pi_{\unsure{0|00}}(1)=\ell_1$)
